% Appendix --- reproducibility details
% Compiled inline with main.tex via % Appendix --- reproducibility details
% Compiled inline with main.tex via % Appendix --- reproducibility details
% Compiled inline with main.tex via % Appendix --- reproducibility details
% Compiled inline with main.tex via \input{appendix.tex} after \appendix

\section{Reproducibility}
\label{sec:reproducibility}

\subsection{Code and data}
Code public at \url{https://github.com/vudang4494/xling-grpo-sub3b} (Apache-2.0).
Three LoRA adapters (one per arm), all evaluation JSONs (full \texttt{responses[]}
arrays included), per-step training logs, and reward time series will be released
with the paper.

\subsection{Hyperparameters}
Full configs in \texttt{configs/}; summary below.

\paragraph{GRPO (Open-RS RS2 verbatim, except LoRA fallback).}
$\text{lr}=10^{-6}$, $G{=}6$, \texttt{max\_completion\_length}=3584,
\texttt{max\_prompt\_length}=512, $T_{\text{rollout}}=0.7$, KL $\beta{=}0.04$, 50 steps,
save every 50 steps, scheduler \texttt{cosine\_with\_min\_lr}
(\texttt{min\_lr\_rate}=0.1, \texttt{warmup\_ratio}=0.1), bf16, flash-attention 2,
gradient checkpointing. Effective batch 96 prompts (per-device 6 $\times$
\texttt{gradient\_accumulation\_steps}=16) --- Open-RS used 4$\times$A40 (per-device 6
$\times$ grad\_accum 4 $\times$ 4 GPUs); we use 1$\times$A100 with grad\_accum 16
to match.

\paragraph{LoRA configuration.} $r{=}16$, $\alpha{=}32$, dropout 0.05,
\texttt{bias=none}, \texttt{task\_type=CAUSAL\_LM}, \texttt{target\_modules}
$=\{q, k, v, o\}$\_proj. Open-RS used full-parameter; LoRA $r{=}16$ is our
single-A100 80GB fallback to fit the same effective batch size and rollout
parallelism.

\paragraph{vLLM (in-process rollout).} TRL 0.15.2 spawns vLLM 0.7.2 within the
training process. \texttt{gpu\_memory\_utilization}=0.25 (20 GB), leaving
$\sim$60 GB for training. \texttt{vllm\_max\_model\_len}=4608.

\paragraph{Reward weights.} Arms \armA{} and \armB{}: $w_{R_1}{=}w_{R_2}{=}1.0$.
Arm \armC{}: $w_{R_1}{=}w_{R_2}{=}w_{R_5}{=}1.0$ with
$R_5$'s \texttt{no\_penalty\_for\_en} flag set to \texttt{False} so $R_5$ fires on
every prompt (not just non-EN).

\paragraph{Eval.} vLLM 0.7.2, dtype bfloat16, $T{=}0$ for pass@1 (greedy), $T{=}0.7$
top\_p=0.95 for AIME-2024 maj@8, max\_tokens=8192, no early stop tokens.
Open-RS \texttt{lighteval} \texttt{MATH\_QUERY\_TEMPLATE} verbatim; no system
prompt (template embedded in user message).

\subsection{Datasets (HuggingFace IDs)}
\begin{itemize}\itemsep0pt
\item \armA{}, \armC{} training: \texttt{knoveleng/open-rs} (MIT, 7K, \texttt{split=train}).
\item \armB{} training:
\path{5CD-AI/Vietnamese-meta-math-MetaMathQA-40K-gg-translated}; license unset
on Hub (parent MetaMathQA = MIT). We extracted 5{,}203 records from a 7K random
subset by retaining only those whose Vietnamese response contains the
``\DJ\'ap \'an l\`a'' (Vietnamese for ``the answer is'') answer marker.
\item Eval: \path{HuggingFaceH4/MATH-500} (500 test), \path{Maxwell-Jia/AIME_2024}
(MIT, 30 records; field names capitalized: \texttt{Problem}, \texttt{Solution},
\texttt{Answer}, \texttt{ID}), \path{knoveleng/AMC-23} (40 records; license
unspecified on Hub, originally AoPS wiki content).
\end{itemize}

\subsection{Decontamination}
We did not perform train/test decontamination for the present arms because the
training datasets (\texttt{knoveleng/open-rs} and \texttt{5CD-AI/Vietnamese-MetaMathQA})
are upstream of the evaluation datasets in different domains; in addition,
Open-RS RS2 itself does not decontaminate. A future ablation will quantify any
contamination effect via 8-gram overlap statistics.

\subsection{Compute resources}
Training and evaluation: 1$\times$ NVIDIA A100-SXM4-80GB GPU. Total
wallclock for the three arms (50 steps each
plus ckpt-50 evaluation on three benchmarks): 11 hours 39 minutes.

\subsection{Random seeds}
Each cell (arm, seed) uses one of \texttt{seed} $\in \{42, 123, 7\}$ for
GRPO training. AIME-2024 \texttt{maj@8} uses seeds 42--49 (8 stochastic
samples per problem). AMC-23 \texttt{maj@4} uses seeds 42--45. Pass@1
evaluations are deterministic ($T{=}0$) and seed-independent. The
$\base$/\armD{} AMC-23 \texttt{pass@1} cells are single-seed
(\texttt{42}); all other cells report three seeds.

\subsection{Software versions}
Python 3.12.3, \texttt{torch==2.5.1+cu124}, \texttt{transformers==4.49.0},
\texttt{trl==0.15.2}, \texttt{tokenizers==0.21.0}, \texttt{vllm==0.7.2},
\texttt{accelerate==1.2.1}, \texttt{peft==0.14.0}, \texttt{datasets==4.8.5},
\texttt{flash\_attn==2.7.2.post1}, \texttt{sympy==1.13.1},
\texttt{math\_verify==0.9.0}, \texttt{fasttext-wheel==0.9.2}
(with \texttt{lid.176.bin} 126 MB).

\textbf{Note on TRL version.} \texttt{GRPOTrainer} first appeared in TRL 0.14.0;
\texttt{GRPOConfig.reward\_weights} field requires TRL $\geq$ 0.15.0. TRL 0.15.2
spawns vLLM in-process; TRL 0.16.0+ requires an external vLLM server. Our pipeline
targets the 0.15.2 in-process model.

\textbf{Note on sympy version.} Earlier project notes pinned
\texttt{sympy<1.13} citing potential Math-Verify breakage; in practice
\texttt{math\_verify==0.9.0} works correctly with \texttt{sympy==1.13.1}
(verified by running \texttt{verify(parse(``1/2''), parse(``0.5''))}).

\subsection{Hardware notes}
The single-A100 LoRA constraint required several config adjustments not present
in Open-RS:
\begin{itemize}\itemsep0pt
\item \texttt{vllm\_gpu\_memory\_utilization} reduced from Open-RS's 0.7 to 0.25 to
leave room for the LoRA training pass.
\item Effective batch maintained at 96 (Open-RS) by setting
\texttt{gradient\_accumulation\_steps}=16 instead of 4 (4$\times$A40) on a single
device.
\item \texttt{gradient\_checkpointing} enabled with
\texttt{use\_reentrant=False} to fit activations in remaining VRAM.
\end{itemize}

\subsection{Limitations not in main text}
\begin{itemize}\itemsep0pt
\item \textbf{Evaluation gap to Open-RS at base.} Even with verbatim Open-RS
\texttt{lighteval} prompt, our base \texttt{pass@1} differs from
Open-RS-reported numbers by $-12.9$\,pp on AMC-23 (50.0\,\% vs 62.9\,\%) and
$-23.4$\,pp on MATH-500 (59.4\,\% vs 82.8\,\%). The gap closes under
like-for-like \texttt{maj@k} comparison
(Section~\ref{sec:limitations}, metric-mismatch discussion). Suspected residual causes: Math-Verify configuration differences
between our adapter and Open-RS's parser (boxed-match priority), or vLLM
tokenizer edge cases. We report numbers vs.\ our own base rather than vs.\
Open-RS reported.
\item \textbf{AIME-2024 baseline closer to Open-RS.} $26.7$\,\% vs $28.8$\,\%,
gap $-2.1$\,pp --- suggesting evaluation pipeline issues are
problem-difficulty-dependent.
\end{itemize}

% =============================================================================
\section{Data Provenance}
\label{sec:data_provenance}

Every number in Tables~\ref{tab:main} and~\ref{tab:delta} traces to an
evaluation JSON file in \path{results/eval/}.  Table~\ref{tab:sources}
gives the exact source file for each data point.  Two distinct base models are
used across benchmarks because different evaluation batches were run at different
times:
\begin{itemize}\itemsep0pt
\item \texttt{base\_v3\_openrs\_eval/}: checkpointed at training step 0 with the
same pipeline used for the trained arms; provides MATH-500 and AIME-2024
benchmarks for the base.
\item \texttt{w17\_base\_distill\_v2/}: raw pretrained
\path{deepseek-ai/DeepSeek-R1-Distill-Qwen-1.5B} re-evaluated in the W17 batch
with \texttt{temperature\_maj}=0.7; provides the AMC-23 \texttt{maj@4} baseline
because the v3 eval was affected by the temperature bug.
\item \texttt{base\_distill15b\_mgsm/}: raw pretrained model used for MGSM
evaluation only (a separate eval batch).
\end{itemize}

\begin{table*}[t]
\centering\footnotesize
\caption{Source JSON files for every number in Tables~\ref{tab:main} and~\ref{tab:delta}.
All paths are relative to \texttt{results/eval/}.}
\label{tab:sources}
\begin{tabular}{@{}p{0.30\textwidth} l p{0.46\textwidth}@{}}
\toprule
Cell & Metric & Source JSON file \\
\midrule
\base{} AMC-23 \texttt{maj@4} & $70.0\%$ &
\path{w17\_base\_distill\_v2/base\_distill\_v2\_amc23.json} \\
\base{} all other metrics & Table~\ref{tab:main} &
\path{base\_v3\_openrs\_eval/base\_v3\_\{bench\}.json} \\
\armA{} seeds 42/123/7 AMC-23 \texttt{maj@4} &
$70.0\pm4.3\%$ &
\path{w17\_reproduce\_openrs\_rs2\_\{seed\}\_v2/...\_amc23.json} \\
\armB{} seeds 42/123/7 AMC-23 \texttt{maj@4} &
$70.0\pm2.5\%$ &
\path{w17\_a2\_vi\_\{seed\}\_v2/...\_amc23.json} \\
\armC{} seeds 42/123/7 AMC-23 \texttt{maj@4} &
$68.3\pm3.8\%$ &
\path{w17\_a3\_enlang\_\{seed\}\_v2/...\_amc23.json} \\
\armD{} all metrics &
Table~\ref{tab:main} &
\path{a4\_const\_bias\_\{seed\}\_step50/...\_\{bench\}.json} \\
\midrule
\multicolumn{3}{l}{All other cells / benchmarks / metrics} \\
\phantom{\armA{}} & &
\path{\{cell\}\_step50/...\_\{bench\}.json} \\
\bottomrule
\end{tabular}
\end{table*}

\subsection{AMC-23 \texttt{maj@4}: temperature bug and fix}
\label{sec:maj4_bug}
During initial evaluation (May 2026) the \texttt{configs/eval.yaml} was missing
the \texttt{temperature\_maj} field.  Consequently all majority-vote rollouts used
\texttt{temperature}=0.0 (greedy), producing \emph{identical} samples for every
seed.  Majority vote collapsed to a single prediction; \texttt{maj@4}=0 for all
cells using seed 42 (the majority-vote seeds start at 42, 43, 44, 45).  The
bug was diagnosed on May 15 (W17) and fixed by adding \texttt{temperature\_maj:
0.7} to the config; all trained-arm AMC-23 evals were re-run as
\path{w17\_*\_v2/} directories.

The \base{} model was not re-evaluated; its reported \texttt{maj@4}=70.0\%
comes from \path{w17\_base\_distill\_v2/} (the W17 base batch that used the
correct temperature).  The \armD{} arm was not included in the W17 re-run
and retains its original step-50 values; these are already valid because
their seed-42 majority-vote seeds (42, 43, 44, 45) happened to use
temperature 0.7 by accident in the original run (the fallback in
\texttt{src/eval/runner.py} defaults to 0.7 when
\texttt{temperature\_maj} is absent), so no correction was needed.

\subsection{Training-curve logging limitation}
The released training-curve logs report mean $\pm \sigma$ across only
\emph{two} seeds (123 and 7) for
arms \armA{}, \armB{}, \armC{}, because the seed-42 checkpoint logs were not
saved for those arms.  The \armD{} training curves are not included because
its training was run separately after the main logging infrastructure was
updated.  This does not affect the main paper claims, which are based on
evaluation results across all three seeds.
 after \appendix

\section{Reproducibility}
\label{sec:reproducibility}

\subsection{Code and data}
Code public at \url{https://github.com/vudang4494/xling-grpo-sub3b} (Apache-2.0).
Three LoRA adapters (one per arm), all evaluation JSONs (full \texttt{responses[]}
arrays included), per-step training logs, and reward time series will be released
with the paper.

\subsection{Hyperparameters}
Full configs in \texttt{configs/}; summary below.

\paragraph{GRPO (Open-RS RS2 verbatim, except LoRA fallback).}
$\text{lr}=10^{-6}$, $G{=}6$, \texttt{max\_completion\_length}=3584,
\texttt{max\_prompt\_length}=512, $T_{\text{rollout}}=0.7$, KL $\beta{=}0.04$, 50 steps,
save every 50 steps, scheduler \texttt{cosine\_with\_min\_lr}
(\texttt{min\_lr\_rate}=0.1, \texttt{warmup\_ratio}=0.1), bf16, flash-attention 2,
gradient checkpointing. Effective batch 96 prompts (per-device 6 $\times$
\texttt{gradient\_accumulation\_steps}=16) --- Open-RS used 4$\times$A40 (per-device 6
$\times$ grad\_accum 4 $\times$ 4 GPUs); we use 1$\times$A100 with grad\_accum 16
to match.

\paragraph{LoRA configuration.} $r{=}16$, $\alpha{=}32$, dropout 0.05,
\texttt{bias=none}, \texttt{task\_type=CAUSAL\_LM}, \texttt{target\_modules}
$=\{q, k, v, o\}$\_proj. Open-RS used full-parameter; LoRA $r{=}16$ is our
single-A100 80GB fallback to fit the same effective batch size and rollout
parallelism.

\paragraph{vLLM (in-process rollout).} TRL 0.15.2 spawns vLLM 0.7.2 within the
training process. \texttt{gpu\_memory\_utilization}=0.25 (20 GB), leaving
$\sim$60 GB for training. \texttt{vllm\_max\_model\_len}=4608.

\paragraph{Reward weights.} Arms \armA{} and \armB{}: $w_{R_1}{=}w_{R_2}{=}1.0$.
Arm \armC{}: $w_{R_1}{=}w_{R_2}{=}w_{R_5}{=}1.0$ with
$R_5$'s \texttt{no\_penalty\_for\_en} flag set to \texttt{False} so $R_5$ fires on
every prompt (not just non-EN).

\paragraph{Eval.} vLLM 0.7.2, dtype bfloat16, $T{=}0$ for pass@1 (greedy), $T{=}0.7$
top\_p=0.95 for AIME-2024 maj@8, max\_tokens=8192, no early stop tokens.
Open-RS \texttt{lighteval} \texttt{MATH\_QUERY\_TEMPLATE} verbatim; no system
prompt (template embedded in user message).

\subsection{Datasets (HuggingFace IDs)}
\begin{itemize}\itemsep0pt
\item \armA{}, \armC{} training: \texttt{knoveleng/open-rs} (MIT, 7K, \texttt{split=train}).
\item \armB{} training:
\path{5CD-AI/Vietnamese-meta-math-MetaMathQA-40K-gg-translated}; license unset
on Hub (parent MetaMathQA = MIT). We extracted 5{,}203 records from a 7K random
subset by retaining only those whose Vietnamese response contains the
``\DJ\'ap \'an l\`a'' (Vietnamese for ``the answer is'') answer marker.
\item Eval: \path{HuggingFaceH4/MATH-500} (500 test), \path{Maxwell-Jia/AIME_2024}
(MIT, 30 records; field names capitalized: \texttt{Problem}, \texttt{Solution},
\texttt{Answer}, \texttt{ID}), \path{knoveleng/AMC-23} (40 records; license
unspecified on Hub, originally AoPS wiki content).
\end{itemize}

\subsection{Decontamination}
We did not perform train/test decontamination for the present arms because the
training datasets (\texttt{knoveleng/open-rs} and \texttt{5CD-AI/Vietnamese-MetaMathQA})
are upstream of the evaluation datasets in different domains; in addition,
Open-RS RS2 itself does not decontaminate. A future ablation will quantify any
contamination effect via 8-gram overlap statistics.

\subsection{Compute resources}
Training and evaluation: 1$\times$ NVIDIA A100-SXM4-80GB GPU. Total
wallclock for the three arms (50 steps each
plus ckpt-50 evaluation on three benchmarks): 11 hours 39 minutes.

\subsection{Random seeds}
Each cell (arm, seed) uses one of \texttt{seed} $\in \{42, 123, 7\}$ for
GRPO training. AIME-2024 \texttt{maj@8} uses seeds 42--49 (8 stochastic
samples per problem). AMC-23 \texttt{maj@4} uses seeds 42--45. Pass@1
evaluations are deterministic ($T{=}0$) and seed-independent. The
$\base$/\armD{} AMC-23 \texttt{pass@1} cells are single-seed
(\texttt{42}); all other cells report three seeds.

\subsection{Software versions}
Python 3.12.3, \texttt{torch==2.5.1+cu124}, \texttt{transformers==4.49.0},
\texttt{trl==0.15.2}, \texttt{tokenizers==0.21.0}, \texttt{vllm==0.7.2},
\texttt{accelerate==1.2.1}, \texttt{peft==0.14.0}, \texttt{datasets==4.8.5},
\texttt{flash\_attn==2.7.2.post1}, \texttt{sympy==1.13.1},
\texttt{math\_verify==0.9.0}, \texttt{fasttext-wheel==0.9.2}
(with \texttt{lid.176.bin} 126 MB).

\textbf{Note on TRL version.} \texttt{GRPOTrainer} first appeared in TRL 0.14.0;
\texttt{GRPOConfig.reward\_weights} field requires TRL $\geq$ 0.15.0. TRL 0.15.2
spawns vLLM in-process; TRL 0.16.0+ requires an external vLLM server. Our pipeline
targets the 0.15.2 in-process model.

\textbf{Note on sympy version.} Earlier project notes pinned
\texttt{sympy<1.13} citing potential Math-Verify breakage; in practice
\texttt{math\_verify==0.9.0} works correctly with \texttt{sympy==1.13.1}
(verified by running \texttt{verify(parse(``1/2''), parse(``0.5''))}).

\subsection{Hardware notes}
The single-A100 LoRA constraint required several config adjustments not present
in Open-RS:
\begin{itemize}\itemsep0pt
\item \texttt{vllm\_gpu\_memory\_utilization} reduced from Open-RS's 0.7 to 0.25 to
leave room for the LoRA training pass.
\item Effective batch maintained at 96 (Open-RS) by setting
\texttt{gradient\_accumulation\_steps}=16 instead of 4 (4$\times$A40) on a single
device.
\item \texttt{gradient\_checkpointing} enabled with
\texttt{use\_reentrant=False} to fit activations in remaining VRAM.
\end{itemize}

\subsection{Limitations not in main text}
\begin{itemize}\itemsep0pt
\item \textbf{Evaluation gap to Open-RS at base.} Even with verbatim Open-RS
\texttt{lighteval} prompt, our base \texttt{pass@1} differs from
Open-RS-reported numbers by $-12.9$\,pp on AMC-23 (50.0\,\% vs 62.9\,\%) and
$-23.4$\,pp on MATH-500 (59.4\,\% vs 82.8\,\%). The gap closes under
like-for-like \texttt{maj@k} comparison
(Section~\ref{sec:limitations}, metric-mismatch discussion). Suspected residual causes: Math-Verify configuration differences
between our adapter and Open-RS's parser (boxed-match priority), or vLLM
tokenizer edge cases. We report numbers vs.\ our own base rather than vs.\
Open-RS reported.
\item \textbf{AIME-2024 baseline closer to Open-RS.} $26.7$\,\% vs $28.8$\,\%,
gap $-2.1$\,pp --- suggesting evaluation pipeline issues are
problem-difficulty-dependent.
\end{itemize}

% =============================================================================
\section{Data Provenance}
\label{sec:data_provenance}

Every number in Tables~\ref{tab:main} and~\ref{tab:delta} traces to an
evaluation JSON file in \path{results/eval/}.  Table~\ref{tab:sources}
gives the exact source file for each data point.  Two distinct base models are
used across benchmarks because different evaluation batches were run at different
times:
\begin{itemize}\itemsep0pt
\item \texttt{base\_v3\_openrs\_eval/}: checkpointed at training step 0 with the
same pipeline used for the trained arms; provides MATH-500 and AIME-2024
benchmarks for the base.
\item \texttt{w17\_base\_distill\_v2/}: raw pretrained
\path{deepseek-ai/DeepSeek-R1-Distill-Qwen-1.5B} re-evaluated in the W17 batch
with \texttt{temperature\_maj}=0.7; provides the AMC-23 \texttt{maj@4} baseline
because the v3 eval was affected by the temperature bug.
\item \texttt{base\_distill15b\_mgsm/}: raw pretrained model used for MGSM
evaluation only (a separate eval batch).
\end{itemize}

\begin{table*}[t]
\centering\footnotesize
\caption{Source JSON files for every number in Tables~\ref{tab:main} and~\ref{tab:delta}.
All paths are relative to \texttt{results/eval/}.}
\label{tab:sources}
\begin{tabular}{@{}p{0.30\textwidth} l p{0.46\textwidth}@{}}
\toprule
Cell & Metric & Source JSON file \\
\midrule
\base{} AMC-23 \texttt{maj@4} & $70.0\%$ &
\path{w17\_base\_distill\_v2/base\_distill\_v2\_amc23.json} \\
\base{} all other metrics & Table~\ref{tab:main} &
\path{base\_v3\_openrs\_eval/base\_v3\_\{bench\}.json} \\
\armA{} seeds 42/123/7 AMC-23 \texttt{maj@4} &
$70.0\pm4.3\%$ &
\path{w17\_reproduce\_openrs\_rs2\_\{seed\}\_v2/...\_amc23.json} \\
\armB{} seeds 42/123/7 AMC-23 \texttt{maj@4} &
$70.0\pm2.5\%$ &
\path{w17\_a2\_vi\_\{seed\}\_v2/...\_amc23.json} \\
\armC{} seeds 42/123/7 AMC-23 \texttt{maj@4} &
$68.3\pm3.8\%$ &
\path{w17\_a3\_enlang\_\{seed\}\_v2/...\_amc23.json} \\
\armD{} all metrics &
Table~\ref{tab:main} &
\path{a4\_const\_bias\_\{seed\}\_step50/...\_\{bench\}.json} \\
\midrule
\multicolumn{3}{l}{All other cells / benchmarks / metrics} \\
\phantom{\armA{}} & &
\path{\{cell\}\_step50/...\_\{bench\}.json} \\
\bottomrule
\end{tabular}
\end{table*}

\subsection{AMC-23 \texttt{maj@4}: temperature bug and fix}
\label{sec:maj4_bug}
During initial evaluation (May 2026) the \texttt{configs/eval.yaml} was missing
the \texttt{temperature\_maj} field.  Consequently all majority-vote rollouts used
\texttt{temperature}=0.0 (greedy), producing \emph{identical} samples for every
seed.  Majority vote collapsed to a single prediction; \texttt{maj@4}=0 for all
cells using seed 42 (the majority-vote seeds start at 42, 43, 44, 45).  The
bug was diagnosed on May 15 (W17) and fixed by adding \texttt{temperature\_maj:
0.7} to the config; all trained-arm AMC-23 evals were re-run as
\path{w17\_*\_v2/} directories.

The \base{} model was not re-evaluated; its reported \texttt{maj@4}=70.0\%
comes from \path{w17\_base\_distill\_v2/} (the W17 base batch that used the
correct temperature).  The \armD{} arm was not included in the W17 re-run
and retains its original step-50 values; these are already valid because
their seed-42 majority-vote seeds (42, 43, 44, 45) happened to use
temperature 0.7 by accident in the original run (the fallback in
\texttt{src/eval/runner.py} defaults to 0.7 when
\texttt{temperature\_maj} is absent), so no correction was needed.

\subsection{Training-curve logging limitation}
The released training-curve logs report mean $\pm \sigma$ across only
\emph{two} seeds (123 and 7) for
arms \armA{}, \armB{}, \armC{}, because the seed-42 checkpoint logs were not
saved for those arms.  The \armD{} training curves are not included because
its training was run separately after the main logging infrastructure was
updated.  This does not affect the main paper claims, which are based on
evaluation results across all three seeds.
 after \appendix

\section{Reproducibility}
\label{sec:reproducibility}

\subsection{Code and data}
Code public at \url{https://github.com/vudang4494/xling-grpo-sub3b} (Apache-2.0).
Three LoRA adapters (one per arm), all evaluation JSONs (full \texttt{responses[]}
arrays included), per-step training logs, and reward time series will be released
with the paper.

\subsection{Hyperparameters}
Full configs in \texttt{configs/}; summary below.

\paragraph{GRPO (Open-RS RS2 verbatim, except LoRA fallback).}
$\text{lr}=10^{-6}$, $G{=}6$, \texttt{max\_completion\_length}=3584,
\texttt{max\_prompt\_length}=512, $T_{\text{rollout}}=0.7$, KL $\beta{=}0.04$, 50 steps,
save every 50 steps, scheduler \texttt{cosine\_with\_min\_lr}
(\texttt{min\_lr\_rate}=0.1, \texttt{warmup\_ratio}=0.1), bf16, flash-attention 2,
gradient checkpointing. Effective batch 96 prompts (per-device 6 $\times$
\texttt{gradient\_accumulation\_steps}=16) --- Open-RS used 4$\times$A40 (per-device 6
$\times$ grad\_accum 4 $\times$ 4 GPUs); we use 1$\times$A100 with grad\_accum 16
to match.

\paragraph{LoRA configuration.} $r{=}16$, $\alpha{=}32$, dropout 0.05,
\texttt{bias=none}, \texttt{task\_type=CAUSAL\_LM}, \texttt{target\_modules}
$=\{q, k, v, o\}$\_proj. Open-RS used full-parameter; LoRA $r{=}16$ is our
single-A100 80GB fallback to fit the same effective batch size and rollout
parallelism.

\paragraph{vLLM (in-process rollout).} TRL 0.15.2 spawns vLLM 0.7.2 within the
training process. \texttt{gpu\_memory\_utilization}=0.25 (20 GB), leaving
$\sim$60 GB for training. \texttt{vllm\_max\_model\_len}=4608.

\paragraph{Reward weights.} Arms \armA{} and \armB{}: $w_{R_1}{=}w_{R_2}{=}1.0$.
Arm \armC{}: $w_{R_1}{=}w_{R_2}{=}w_{R_5}{=}1.0$ with
$R_5$'s \texttt{no\_penalty\_for\_en} flag set to \texttt{False} so $R_5$ fires on
every prompt (not just non-EN).

\paragraph{Eval.} vLLM 0.7.2, dtype bfloat16, $T{=}0$ for pass@1 (greedy), $T{=}0.7$
top\_p=0.95 for AIME-2024 maj@8, max\_tokens=8192, no early stop tokens.
Open-RS \texttt{lighteval} \texttt{MATH\_QUERY\_TEMPLATE} verbatim; no system
prompt (template embedded in user message).

\subsection{Datasets (HuggingFace IDs)}
\begin{itemize}\itemsep0pt
\item \armA{}, \armC{} training: \texttt{knoveleng/open-rs} (MIT, 7K, \texttt{split=train}).
\item \armB{} training:
\path{5CD-AI/Vietnamese-meta-math-MetaMathQA-40K-gg-translated}; license unset
on Hub (parent MetaMathQA = MIT). We extracted 5{,}203 records from a 7K random
subset by retaining only those whose Vietnamese response contains the
``\DJ\'ap \'an l\`a'' (Vietnamese for ``the answer is'') answer marker.
\item Eval: \path{HuggingFaceH4/MATH-500} (500 test), \path{Maxwell-Jia/AIME_2024}
(MIT, 30 records; field names capitalized: \texttt{Problem}, \texttt{Solution},
\texttt{Answer}, \texttt{ID}), \path{knoveleng/AMC-23} (40 records; license
unspecified on Hub, originally AoPS wiki content).
\end{itemize}

\subsection{Decontamination}
We did not perform train/test decontamination for the present arms because the
training datasets (\texttt{knoveleng/open-rs} and \texttt{5CD-AI/Vietnamese-MetaMathQA})
are upstream of the evaluation datasets in different domains; in addition,
Open-RS RS2 itself does not decontaminate. A future ablation will quantify any
contamination effect via 8-gram overlap statistics.

\subsection{Compute resources}
Training and evaluation: 1$\times$ NVIDIA A100-SXM4-80GB GPU. Total
wallclock for the three arms (50 steps each
plus ckpt-50 evaluation on three benchmarks): 11 hours 39 minutes.

\subsection{Random seeds}
Each cell (arm, seed) uses one of \texttt{seed} $\in \{42, 123, 7\}$ for
GRPO training. AIME-2024 \texttt{maj@8} uses seeds 42--49 (8 stochastic
samples per problem). AMC-23 \texttt{maj@4} uses seeds 42--45. Pass@1
evaluations are deterministic ($T{=}0$) and seed-independent. The
$\base$/\armD{} AMC-23 \texttt{pass@1} cells are single-seed
(\texttt{42}); all other cells report three seeds.

\subsection{Software versions}
Python 3.12.3, \texttt{torch==2.5.1+cu124}, \texttt{transformers==4.49.0},
\texttt{trl==0.15.2}, \texttt{tokenizers==0.21.0}, \texttt{vllm==0.7.2},
\texttt{accelerate==1.2.1}, \texttt{peft==0.14.0}, \texttt{datasets==4.8.5},
\texttt{flash\_attn==2.7.2.post1}, \texttt{sympy==1.13.1},
\texttt{math\_verify==0.9.0}, \texttt{fasttext-wheel==0.9.2}
(with \texttt{lid.176.bin} 126 MB).

\textbf{Note on TRL version.} \texttt{GRPOTrainer} first appeared in TRL 0.14.0;
\texttt{GRPOConfig.reward\_weights} field requires TRL $\geq$ 0.15.0. TRL 0.15.2
spawns vLLM in-process; TRL 0.16.0+ requires an external vLLM server. Our pipeline
targets the 0.15.2 in-process model.

\textbf{Note on sympy version.} Earlier project notes pinned
\texttt{sympy<1.13} citing potential Math-Verify breakage; in practice
\texttt{math\_verify==0.9.0} works correctly with \texttt{sympy==1.13.1}
(verified by running \texttt{verify(parse(``1/2''), parse(``0.5''))}).

\subsection{Hardware notes}
The single-A100 LoRA constraint required several config adjustments not present
in Open-RS:
\begin{itemize}\itemsep0pt
\item \texttt{vllm\_gpu\_memory\_utilization} reduced from Open-RS's 0.7 to 0.25 to
leave room for the LoRA training pass.
\item Effective batch maintained at 96 (Open-RS) by setting
\texttt{gradient\_accumulation\_steps}=16 instead of 4 (4$\times$A40) on a single
device.
\item \texttt{gradient\_checkpointing} enabled with
\texttt{use\_reentrant=False} to fit activations in remaining VRAM.
\end{itemize}

\subsection{Limitations not in main text}
\begin{itemize}\itemsep0pt
\item \textbf{Evaluation gap to Open-RS at base.} Even with verbatim Open-RS
\texttt{lighteval} prompt, our base \texttt{pass@1} differs from
Open-RS-reported numbers by $-12.9$\,pp on AMC-23 (50.0\,\% vs 62.9\,\%) and
$-23.4$\,pp on MATH-500 (59.4\,\% vs 82.8\,\%). The gap closes under
like-for-like \texttt{maj@k} comparison
(Section~\ref{sec:limitations}, metric-mismatch discussion). Suspected residual causes: Math-Verify configuration differences
between our adapter and Open-RS's parser (boxed-match priority), or vLLM
tokenizer edge cases. We report numbers vs.\ our own base rather than vs.\
Open-RS reported.
\item \textbf{AIME-2024 baseline closer to Open-RS.} $26.7$\,\% vs $28.8$\,\%,
gap $-2.1$\,pp --- suggesting evaluation pipeline issues are
problem-difficulty-dependent.
\end{itemize}

% =============================================================================
\section{Data Provenance}
\label{sec:data_provenance}

Every number in Tables~\ref{tab:main} and~\ref{tab:delta} traces to an
evaluation JSON file in \path{results/eval/}.  Table~\ref{tab:sources}
gives the exact source file for each data point.  Two distinct base models are
used across benchmarks because different evaluation batches were run at different
times:
\begin{itemize}\itemsep0pt
\item \texttt{base\_v3\_openrs\_eval/}: checkpointed at training step 0 with the
same pipeline used for the trained arms; provides MATH-500 and AIME-2024
benchmarks for the base.
\item \texttt{w17\_base\_distill\_v2/}: raw pretrained
\path{deepseek-ai/DeepSeek-R1-Distill-Qwen-1.5B} re-evaluated in the W17 batch
with \texttt{temperature\_maj}=0.7; provides the AMC-23 \texttt{maj@4} baseline
because the v3 eval was affected by the temperature bug.
\item \texttt{base\_distill15b\_mgsm/}: raw pretrained model used for MGSM
evaluation only (a separate eval batch).
\end{itemize}

\begin{table*}[t]
\centering\footnotesize
\caption{Source JSON files for every number in Tables~\ref{tab:main} and~\ref{tab:delta}.
All paths are relative to \texttt{results/eval/}.}
\label{tab:sources}
\begin{tabular}{@{}p{0.30\textwidth} l p{0.46\textwidth}@{}}
\toprule
Cell & Metric & Source JSON file \\
\midrule
\base{} AMC-23 \texttt{maj@4} & $70.0\%$ &
\path{w17\_base\_distill\_v2/base\_distill\_v2\_amc23.json} \\
\base{} all other metrics & Table~\ref{tab:main} &
\path{base\_v3\_openrs\_eval/base\_v3\_\{bench\}.json} \\
\armA{} seeds 42/123/7 AMC-23 \texttt{maj@4} &
$70.0\pm4.3\%$ &
\path{w17\_reproduce\_openrs\_rs2\_\{seed\}\_v2/...\_amc23.json} \\
\armB{} seeds 42/123/7 AMC-23 \texttt{maj@4} &
$70.0\pm2.5\%$ &
\path{w17\_a2\_vi\_\{seed\}\_v2/...\_amc23.json} \\
\armC{} seeds 42/123/7 AMC-23 \texttt{maj@4} &
$68.3\pm3.8\%$ &
\path{w17\_a3\_enlang\_\{seed\}\_v2/...\_amc23.json} \\
\armD{} all metrics &
Table~\ref{tab:main} &
\path{a4\_const\_bias\_\{seed\}\_step50/...\_\{bench\}.json} \\
\midrule
\multicolumn{3}{l}{All other cells / benchmarks / metrics} \\
\phantom{\armA{}} & &
\path{\{cell\}\_step50/...\_\{bench\}.json} \\
\bottomrule
\end{tabular}
\end{table*}

\subsection{AMC-23 \texttt{maj@4}: temperature bug and fix}
\label{sec:maj4_bug}
During initial evaluation (May 2026) the \texttt{configs/eval.yaml} was missing
the \texttt{temperature\_maj} field.  Consequently all majority-vote rollouts used
\texttt{temperature}=0.0 (greedy), producing \emph{identical} samples for every
seed.  Majority vote collapsed to a single prediction; \texttt{maj@4}=0 for all
cells using seed 42 (the majority-vote seeds start at 42, 43, 44, 45).  The
bug was diagnosed on May 15 (W17) and fixed by adding \texttt{temperature\_maj:
0.7} to the config; all trained-arm AMC-23 evals were re-run as
\path{w17\_*\_v2/} directories.

The \base{} model was not re-evaluated; its reported \texttt{maj@4}=70.0\%
comes from \path{w17\_base\_distill\_v2/} (the W17 base batch that used the
correct temperature).  The \armD{} arm was not included in the W17 re-run
and retains its original step-50 values; these are already valid because
their seed-42 majority-vote seeds (42, 43, 44, 45) happened to use
temperature 0.7 by accident in the original run (the fallback in
\texttt{src/eval/runner.py} defaults to 0.7 when
\texttt{temperature\_maj} is absent), so no correction was needed.

\subsection{Training-curve logging limitation}
The released training-curve logs report mean $\pm \sigma$ across only
\emph{two} seeds (123 and 7) for
arms \armA{}, \armB{}, \armC{}, because the seed-42 checkpoint logs were not
saved for those arms.  The \armD{} training curves are not included because
its training was run separately after the main logging infrastructure was
updated.  This does not affect the main paper claims, which are based on
evaluation results across all three seeds.
 after \appendix

\section{Reproducibility}
\label{sec:reproducibility}

\subsection{Code and data}
Code public at \url{https://github.com/vudang4494/xling-grpo-sub3b} (Apache-2.0).
Three LoRA adapters (one per arm), all evaluation JSONs (full \texttt{responses[]}
arrays included), per-step training logs, and reward time series will be released
with the paper.

\subsection{Hyperparameters}
Full configs in \texttt{configs/}; summary below.

\paragraph{GRPO (Open-RS RS2 verbatim, except LoRA fallback).}
$\text{lr}=10^{-6}$, $G{=}6$, \texttt{max\_completion\_length}=3584,
\texttt{max\_prompt\_length}=512, $T_{\text{rollout}}=0.7$, KL $\beta{=}0.04$, 50 steps,
save every 50 steps, scheduler \texttt{cosine\_with\_min\_lr}
(\texttt{min\_lr\_rate}=0.1, \texttt{warmup\_ratio}=0.1), bf16, flash-attention 2,
gradient checkpointing. Effective batch 96 prompts (per-device 6 $\times$
\texttt{gradient\_accumulation\_steps}=16) --- Open-RS used 4$\times$A40 (per-device 6
$\times$ grad\_accum 4 $\times$ 4 GPUs); we use 1$\times$A100 with grad\_accum 16
to match.

\paragraph{LoRA configuration.} $r{=}16$, $\alpha{=}32$, dropout 0.05,
\texttt{bias=none}, \texttt{task\_type=CAUSAL\_LM}, \texttt{target\_modules}
$=\{q, k, v, o\}$\_proj. Open-RS used full-parameter; LoRA $r{=}16$ is our
single-A100 80GB fallback to fit the same effective batch size and rollout
parallelism.

\paragraph{vLLM (in-process rollout).} TRL 0.15.2 spawns vLLM 0.7.2 within the
training process. \texttt{gpu\_memory\_utilization}=0.25 (20 GB), leaving
$\sim$60 GB for training. \texttt{vllm\_max\_model\_len}=4608.

\paragraph{Reward weights.} Arms \armA{} and \armB{}: $w_{R_1}{=}w_{R_2}{=}1.0$.
Arm \armC{}: $w_{R_1}{=}w_{R_2}{=}w_{R_5}{=}1.0$ with
$R_5$'s \texttt{no\_penalty\_for\_en} flag set to \texttt{False} so $R_5$ fires on
every prompt (not just non-EN).

\paragraph{Eval.} vLLM 0.7.2, dtype bfloat16, $T{=}0$ for pass@1 (greedy), $T{=}0.7$
top\_p=0.95 for AIME-2024 maj@8, max\_tokens=8192, no early stop tokens.
Open-RS \texttt{lighteval} \texttt{MATH\_QUERY\_TEMPLATE} verbatim; no system
prompt (template embedded in user message).

\subsection{Datasets (HuggingFace IDs)}
\begin{itemize}\itemsep0pt
\item \armA{}, \armC{} training: \texttt{knoveleng/open-rs} (MIT, 7K, \texttt{split=train}).
\item \armB{} training:
\path{5CD-AI/Vietnamese-meta-math-MetaMathQA-40K-gg-translated}; license unset
on Hub (parent MetaMathQA = MIT). We extracted 5{,}203 records from a 7K random
subset by retaining only those whose Vietnamese response contains the
``\DJ\'ap \'an l\`a'' (Vietnamese for ``the answer is'') answer marker.
\item Eval: \path{HuggingFaceH4/MATH-500} (500 test), \path{Maxwell-Jia/AIME_2024}
(MIT, 30 records; field names capitalized: \texttt{Problem}, \texttt{Solution},
\texttt{Answer}, \texttt{ID}), \path{knoveleng/AMC-23} (40 records; license
unspecified on Hub, originally AoPS wiki content).
\end{itemize}

\subsection{Decontamination}
We did not perform train/test decontamination for the present arms because the
training datasets (\texttt{knoveleng/open-rs} and \texttt{5CD-AI/Vietnamese-MetaMathQA})
are upstream of the evaluation datasets in different domains; in addition,
Open-RS RS2 itself does not decontaminate. A future ablation will quantify any
contamination effect via 8-gram overlap statistics.

\subsection{Compute resources}
Training and evaluation: 1$\times$ NVIDIA A100-SXM4-80GB GPU. Total
wallclock for the three arms (50 steps each
plus ckpt-50 evaluation on three benchmarks): 11 hours 39 minutes.

\subsection{Random seeds}
Each cell (arm, seed) uses one of \texttt{seed} $\in \{42, 123, 7\}$ for
GRPO training. AIME-2024 \texttt{maj@8} uses seeds 42--49 (8 stochastic
samples per problem). AMC-23 \texttt{maj@4} uses seeds 42--45. Pass@1
evaluations are deterministic ($T{=}0$) and seed-independent. The
$\base$/\armD{} AMC-23 \texttt{pass@1} cells are single-seed
(\texttt{42}); all other cells report three seeds.

\subsection{Software versions}
Python 3.12.3, \texttt{torch==2.5.1+cu124}, \texttt{transformers==4.49.0},
\texttt{trl==0.15.2}, \texttt{tokenizers==0.21.0}, \texttt{vllm==0.7.2},
\texttt{accelerate==1.2.1}, \texttt{peft==0.14.0}, \texttt{datasets==4.8.5},
\texttt{flash\_attn==2.7.2.post1}, \texttt{sympy==1.13.1},
\texttt{math\_verify==0.9.0}, \texttt{fasttext-wheel==0.9.2}
(with \texttt{lid.176.bin} 126 MB).

\textbf{Note on TRL version.} \texttt{GRPOTrainer} first appeared in TRL 0.14.0;
\texttt{GRPOConfig.reward\_weights} field requires TRL $\geq$ 0.15.0. TRL 0.15.2
spawns vLLM in-process; TRL 0.16.0+ requires an external vLLM server. Our pipeline
targets the 0.15.2 in-process model.

\textbf{Note on sympy version.} Earlier project notes pinned
\texttt{sympy<1.13} citing potential Math-Verify breakage; in practice
\texttt{math\_verify==0.9.0} works correctly with \texttt{sympy==1.13.1}
(verified by running \texttt{verify(parse(``1/2''), parse(``0.5''))}).

\subsection{Hardware notes}
The single-A100 LoRA constraint required several config adjustments not present
in Open-RS:
\begin{itemize}\itemsep0pt
\item \texttt{vllm\_gpu\_memory\_utilization} reduced from Open-RS's 0.7 to 0.25 to
leave room for the LoRA training pass.
\item Effective batch maintained at 96 (Open-RS) by setting
\texttt{gradient\_accumulation\_steps}=16 instead of 4 (4$\times$A40) on a single
device.
\item \texttt{gradient\_checkpointing} enabled with
\texttt{use\_reentrant=False} to fit activations in remaining VRAM.
\end{itemize}

\subsection{Limitations not in main text}
\begin{itemize}\itemsep0pt
\item \textbf{Evaluation gap to Open-RS at base.} Even with verbatim Open-RS
\texttt{lighteval} prompt, our base \texttt{pass@1} differs from
Open-RS-reported numbers by $-12.9$\,pp on AMC-23 (50.0\,\% vs 62.9\,\%) and
$-23.4$\,pp on MATH-500 (59.4\,\% vs 82.8\,\%). The gap closes under
like-for-like \texttt{maj@k} comparison
(Section~\ref{sec:limitations}, metric-mismatch discussion). Suspected residual causes: Math-Verify configuration differences
between our adapter and Open-RS's parser (boxed-match priority), or vLLM
tokenizer edge cases. We report numbers vs.\ our own base rather than vs.\
Open-RS reported.
\item \textbf{AIME-2024 baseline closer to Open-RS.} $26.7$\,\% vs $28.8$\,\%,
gap $-2.1$\,pp --- suggesting evaluation pipeline issues are
problem-difficulty-dependent.
\end{itemize}

% =============================================================================
\section{Data Provenance}
\label{sec:data_provenance}

Every number in Tables~\ref{tab:main} and~\ref{tab:delta} traces to an
evaluation JSON file in \path{results/eval/}.  Table~\ref{tab:sources}
gives the exact source file for each data point.  Two distinct base models are
used across benchmarks because different evaluation batches were run at different
times:
\begin{itemize}\itemsep0pt
\item \texttt{base\_v3\_openrs\_eval/}: checkpointed at training step 0 with the
same pipeline used for the trained arms; provides MATH-500 and AIME-2024
benchmarks for the base.
\item \texttt{w17\_base\_distill\_v2/}: raw pretrained
\path{deepseek-ai/DeepSeek-R1-Distill-Qwen-1.5B} re-evaluated in the W17 batch
with \texttt{temperature\_maj}=0.7; provides the AMC-23 \texttt{maj@4} baseline
because the v3 eval was affected by the temperature bug.
\item \texttt{base\_distill15b\_mgsm/}: raw pretrained model used for MGSM
evaluation only (a separate eval batch).
\end{itemize}

\begin{table*}[t]
\centering\footnotesize
\caption{Source JSON files for every number in Tables~\ref{tab:main} and~\ref{tab:delta}.
All paths are relative to \texttt{results/eval/}.}
\label{tab:sources}
\begin{tabular}{@{}p{0.30\textwidth} l p{0.46\textwidth}@{}}
\toprule
Cell & Metric & Source JSON file \\
\midrule
\base{} AMC-23 \texttt{maj@4} & $70.0\%$ &
\path{w17\_base\_distill\_v2/base\_distill\_v2\_amc23.json} \\
\base{} all other metrics & Table~\ref{tab:main} &
\path{base\_v3\_openrs\_eval/base\_v3\_\{bench\}.json} \\
\armA{} seeds 42/123/7 AMC-23 \texttt{maj@4} &
$70.0\pm4.3\%$ &
\path{w17\_reproduce\_openrs\_rs2\_\{seed\}\_v2/...\_amc23.json} \\
\armB{} seeds 42/123/7 AMC-23 \texttt{maj@4} &
$70.0\pm2.5\%$ &
\path{w17\_a2\_vi\_\{seed\}\_v2/...\_amc23.json} \\
\armC{} seeds 42/123/7 AMC-23 \texttt{maj@4} &
$68.3\pm3.8\%$ &
\path{w17\_a3\_enlang\_\{seed\}\_v2/...\_amc23.json} \\
\armD{} all metrics &
Table~\ref{tab:main} &
\path{a4\_const\_bias\_\{seed\}\_step50/...\_\{bench\}.json} \\
\midrule
\multicolumn{3}{l}{All other cells / benchmarks / metrics} \\
\phantom{\armA{}} & &
\path{\{cell\}\_step50/...\_\{bench\}.json} \\
\bottomrule
\end{tabular}
\end{table*}

\subsection{AMC-23 \texttt{maj@4}: temperature bug and fix}
\label{sec:maj4_bug}
During initial evaluation (May 2026) the \texttt{configs/eval.yaml} was missing
the \texttt{temperature\_maj} field.  Consequently all majority-vote rollouts used
\texttt{temperature}=0.0 (greedy), producing \emph{identical} samples for every
seed.  Majority vote collapsed to a single prediction; \texttt{maj@4}=0 for all
cells using seed 42 (the majority-vote seeds start at 42, 43, 44, 45).  The
bug was diagnosed on May 15 (W17) and fixed by adding \texttt{temperature\_maj:
0.7} to the config; all trained-arm AMC-23 evals were re-run as
\path{w17\_*\_v2/} directories.

The \base{} model was not re-evaluated; its reported \texttt{maj@4}=70.0\%
comes from \path{w17\_base\_distill\_v2/} (the W17 base batch that used the
correct temperature).  The \armD{} arm was not included in the W17 re-run
and retains its original step-50 values; these are already valid because
their seed-42 majority-vote seeds (42, 43, 44, 45) happened to use
temperature 0.7 by accident in the original run (the fallback in
\texttt{src/eval/runner.py} defaults to 0.7 when
\texttt{temperature\_maj} is absent), so no correction was needed.

\subsection{Training-curve logging limitation}
The released training-curve logs report mean $\pm \sigma$ across only
\emph{two} seeds (123 and 7) for
arms \armA{}, \armB{}, \armC{}, because the seed-42 checkpoint logs were not
saved for those arms.  The \armD{} training curves are not included because
its training was run separately after the main logging infrastructure was
updated.  This does not affect the main paper claims, which are based on
evaluation results across all three seeds.
